\section{Implementation and Optimizations}
\label{sec:implementation}

This section describes the implementation details of MemSt's core components and the optimization strategies that enable its high performance on long-term memory tasks.

\subsection{System Architecture}

MemSt is implemented as a Rust workspace with multiple crates, each responsible for a specific aspect of the memory system. The architecture follows a layered design with clear separation between storage, retrieval, and interface layers.

\subsubsection{Rust Workspace Structure}

\begin{table}[ht]
\centering
\caption{MemSt workspace crates and their responsibilities.}
\label{tab:crates}
\begin{tabular}{@{}lp{7cm}@{}}
\toprule
\textbf{Crate} & \textbf{Purpose} \\
\midrule
memst-core & Storage engines, search indexes, context assembly \\
memst-repo & Git-like MemRepo with content-addressable storage \\
memst-sleep & Async consolidation pipeline with Tokio \\
memst-extract-v2 & KG extraction with SQLite/DuckDB backends \\
memst-mcp & MCP protocol adapter for Claude/Cursor integration \\
memst-cli & Command-line interface \\
memst-py & Python bindings via PyO3 \\
memst-lib & Shared library and FFI utilities \\
\bottomrule
\end{tabular}
\end{table}

\subsection{Optimized Retrieval Strategies}

Our benchmark evaluation (Section~\ref{sec:evaluation}) identified several high-performance retrieval strategies. We describe their implementations below.

\subsubsection{Hierarchical Summary Retrieval}

The Hierarchical Summary strategy achieves the best overall performance (41.7\% F1) by organizing memory into three levels:

\textbf{Level 1: Session Summaries.} Each conversation session is summarized using extractive techniques. We identify key sentences through:
\begin{itemize}
    \item Named entity density (sentences with $>$2 entities)
    \item Temporal markers (dates, relative times)
    \item Sentiment shifts (indicating important events)
\end{itemize}

The summary for session $s$ is computed as:
\begin{equation}
\text{Summary}(s) = \text{concat}\left(\arg\max_{s_i \in s}^{k} \text{score}(s_i)\right)
\end{equation}
where $\text{score}(s_i) = \alpha \cdot \text{entities}(s_i) + \beta \cdot \text{temporal}(s_i) + \gamma \cdot \text{sentiment\_variance}(s_i)$.

\textbf{Level 2: Key Turns.} Within each session, we identify turns that contain:
\begin{itemize}
    \item User preferences (``I prefer...'', ``I don't like...'')
    \item Action items (``Remember to...'', ``Don't forget...'')
    \item Entity definitions (``X is a...'', ``X works at...'')
\end{itemize}

\textbf{Level 3: Full Context.} The complete dialogue is available for detailed queries, accessed via BM25 or vector search.

\textbf{Dynamic Budget Allocation.} The context assembly allocates the token budget based on question type:
\begin{itemize}
    \item Single-hop: 50\% L2 (key turns), 30\% L1, 20\% L3
    \item Multi-hop: 40\% L1, 40\% L2, 20\% L3
    \item Temporal: 30\% L1, 50\% L2 (temporal turns), 20\% L3
    \item Open-domain: 30\% L1, 30\% L2, 40\% L3
\end{itemize}

\subsubsection{Cascade Retrieval}

The Cascade strategy implements a three-stage filtering pipeline for high-precision retrieval:

\textbf{Stage 1: BM25 Candidate Selection.} Fast sparse retrieval to select top-$k_1$ candidates:
\begin{equation}
\text{BM25}(q, d) = \sum_{t \in q} \text{IDF}(t) \cdot \frac{f(t,d) \cdot (k_1 + 1)}{f(t,d) + k_1 \cdot (1 - b + b \cdot \frac{|d|}{\text{avgdl}})}
\end{equation}
with tuned parameters $k_1 = 0.5$, $b = 0.75$ (determined through ablation studies).

\textbf{Stage 2: Vector Re-ranking.} Dense retrieval using HNSW to refine top-$k_1$ to top-$k_2$:
\begin{equation}
\text{score}(q, d) = \cos(\text{embed}(q), \text{embed}(d))
\end{equation}

\textbf{Stage 3: Cross-encoder Final Ranking.} A lightweight cross-attention model scores the final top-$k_3$ candidates. In our implementation, we use a heuristic-based approximation that considers:
\begin{itemize}
    \item Query-document overlap ratio
    \item Length normalization (prefer medium-length contexts)
    \item Position bias (earlier matches in document)
\end{itemize}

The cascade reduces the candidate space from $N$ documents to $k_3 \ll N$ with progressively more expensive scoring functions, achieving 55.2\% F1 on single-hop questions.

\subsubsection{Temporal Knowledge Graph}

The TemporalKG strategy addresses time-sensitive queries through explicit temporal modeling:

\textbf{Time Expression Extraction.} We extract temporal references using pattern matching:
\begin{itemize}
    \item Absolute dates: ``January 15, 2023''
    \item Relative times: ``last week'', ``two months ago''
    \item Implicit references: ``when I was in college''
\end{itemize}

Each temporal expression $t$ is normalized to a canonical form $t'$ using the conversation timeline as reference.

\textbf{Entity Timeline Construction.} For each entity $e$, we build a timeline:
\begin{equation}
\text{Timeline}(e) = [(t_1, e_1, r_1), (t_2, e_2, r_2), \ldots, (t_n, e_n, r_n)]
\end{equation}
where $t_i$ is the normalized time, $e_i$ is the event description, and $r_i$ is the relation type.

\textbf{Temporal Reasoning.} For queries like ``When did Alice move to SF?'', we:
\begin{enumerate}
    \item Extract entity (``Alice'') and relation (``move to SF'')
    \item Retrieve Alice's timeline: $\text{Timeline}(\text{Alice})$
    \item Match relation against events: $\arg\max_{e_i} \text{sim}(e_i, \text{``move to SF''})$
    \item Return associated time $t_i$
\end{enumerate}

This approach achieves 54.0\% F1 on temporal questions, 120\% better than RAG baseline.

\subsubsection{Multi-hop Entity Graph}

The MultiHopKG strategy handles questions requiring indirect connections (e.g., ``Who is Alice's manager's spouse?'').

\textbf{Relation Extraction.} We extract $(s, r, o)$ triples using dependency parsing patterns:
\begin{itemize}
    \item ``$s$ works at $o$'' $\rightarrow$ (s, works\_at, o)
    \item ``$s$ is married to $o$'' $\rightarrow$ (s, spouse, o)
    \item ``$s$ manages $o$'' $\rightarrow$ (s, manager, o)
\end{itemize}

\textbf{Graph Traversal.} For multi-hop queries, we perform breadth-first search:
\begin{equation}
\text{Answer}(q) = \{o \mid s \xrightarrow{r_1} e_1 \xrightarrow{r_2} \ldots \xrightarrow{r_n} o, n \leq k\}
\end{equation}
where $k$ is the maximum hop count (typically 2-3).

The graph structure enables reasoning across disconnected text segments, achieving 52.8\% F1 on multi-hop questions.

\subsection{Hybrid Search Implementation}

MemSt implements a configurable hybrid search system combining BM25 and HNSW vector search.

\subsubsection{BM25 with Tuned Parameters}

Through ablation studies on LOCOMO, we determined optimal BM25 parameters:
\begin{itemize}
    \item $k_1 = 0.5$ (lower than standard 1.2 for short documents)
    \item $b = 0.75$ (standard length normalization)
\end{itemize}

The lower $k_1$ value reduces saturation for short conversation turns, where term frequency is naturally limited.

\subsubsection{HNSW Vector Index}

We use the HNSW (Hierarchical Navigable Small World) algorithm for approximate nearest neighbor search:
\begin{itemize}
    \item $M = 16$ (maximum neighbors per layer)
    \item $ef\_construction = 200$
    \item $ef\_search = 50$
    \item Dimension: 1024 (OpenAI text-embedding-3-small)
    \item Storage: float16 (50\% memory reduction)
\end{itemize}

Search complexity is $O(\log N)$, enabling scalable retrieval even for very long conversations.

\subsubsection{Reciprocal Rank Fusion}

We combine BM25 and vector rankings using Reciprocal Rank Fusion (RRF):
\begin{equation}
\text{RRF}(d) = \sum_{r \in \{bm25, vec\}} \frac{w_r}{k + \text{rank}_r(d)}
\end{equation}

where $w_r$ is the weight for ranker $r$, and $k = 20$ (tuned constant for stability).

Category-specific weights (determined through validation):
\begin{itemize}
    \item Single-hop: $w_{bm25} = 0.7$, $w_{vec} = 0.3$
    \item Multi-hop: $w_{bm25} = 0.5$, $w_{vec} = 0.5$
    \item Temporal: $w_{bm25} = 0.4$, $w_{vec} = 0.6$
    \item Open-domain: $w_{bm25} = 0.3$, $w_{vec} = 0.7$
\end{itemize}

\subsection{Context Assembly}

The context assembly module constructs the final LLM prompt from retrieved memories, respecting token budgets.

\subsubsection{Token Budget Optimization}

Given token budget $B$, we solve a constrained selection problem:
\begin{align}
\text{maximize} \quad & \sum_{m \in S} R(m, q) \cdot I(m) \\
\text{subject to} \quad & \sum_{m \in S} \text{tokens}(m) \leq B \\
& |S| \leq N_{\max}
\end{align}

where $R(m, q)$ is retrieval relevance and $I(m)$ is memory importance.

We use a greedy algorithm with priority ordering:
\begin{enumerate}
    \item Working memory (always included if exists)
    \item Retrieved memories by relevance score
    \item Recent memories if budget remains
\end{enumerate}

\subsubsection{Memory Deduplication}

To avoid redundant context, we deduplicate memories using:
\begin{itemize}
    \item Exact string match (hash-based)
    \item Semantic similarity threshold ($\cos \text{sim} > 0.95$)
    \item Subsumption check (shorter memory contained in longer)
\end{itemize}

\subsection{Performance Optimizations}

\subsubsection{Indexing Optimizations}

\begin{itemize}
    \item \textbf{Parallel Indexing}: BM25 and HNSW indexes built concurrently using Rayon
    \item \textbf{Incremental Updates}: New memories indexed in batches (every 10 messages)
    \item \textbf{Compressed Storage}: zstd compression for message content (3-5x size reduction)
\end{itemize}

\subsubsection{Query Optimizations}

\begin{itemize}
    \item \textbf{Query Caching}: Embedding vectors cached for repeated queries
    \item \textbf{Prefetching}: Likely next queries predicted and prefetched
    \item \textbf{Early Termination}: Search stops when top-$k$ scores plateau
\end{itemize}

\subsubsection{Memory Tier Management}

Automatic transitions between memory tiers use tuned thresholds:
\begin{itemize}
    \item Working $\rightarrow$ Short-term: After 10 turns or context switch
    \item Short-term $\rightarrow$ Long-term: After 7 days or 100 accesses
    \item Long-term $\rightarrow$ Archival: After 90 days without access
\end{itemize}

\subsection{Python Bindings and API}

MemSt exposes a comprehensive Python API via PyO3:

\begin{verbatim}
import memst

# Initialize store
store = memst.SessionStore("./data")

# Create session with hierarchical indexing
session = store.create_session("My Chat", "gpt-4")

# Add messages (automatically indexed)
store.add_message(session.id, memst.Role.User, "Hello")

# Query with specific strategy
results = store.search(
    query="When did we discuss Python?",
    strategy="temporal_kg",  # or "hierarchical", "cascade", etc.
    limit=10
)
\end{verbatim}

The Rust core ensures performance while Python bindings provide flexibility for research and prototyping.
